% !TeX root = ../papers/nc-p3.tex
\section{異なる二固有値を持つBell直線の全Jordan窓}
\label{sec:split-two-eigenvalue-jordan}
\begin{theorem}\label{thm:split-two-eigenvalue-jordan}
直線型で平方固有値が $1,s^2$, $s\ne0$, $s^2\ne1$、少なくとも一方が非半単純となる
全ての可逆分解窓は、次の表と下記の例外一類で尽きる。全点が標準Koszul左右AS正則で最小周期は1である。
表の正規形の不変量は四文字の量子三角積 $T$ と、指定した無向辺の集合 $E$ である。
\[
\begin{array}{c|c|c}
\text{長さ二の鎖} & a & E\\\hline
1 & k^\times & \{02\}\\
2 & 1,-1 & \{02,03\}\\
2 & s^{-2},-s^{-2} & \{02,12\}
\end{array}
\qquad
 T_{012},T_{013},T_{023},T_{123}= -as,\ -a/s,\ -as,\ -as^3 .
\]
表内・表間の無印同型の必要十分条件は、\emph{一つの}四文字置換が $E,T$ の双方を対応させることである。
三角積は奇置換で逆数にする交代関数として三文字の全順序に延長する。
残る一類は $s^2=-1$ で二鎖のときだけ現れ、閉軌道は片側反転の $a^2=1$ である。
表の閉軌道は全て表示した量子類である。十五次元成分との交差は、表の $s^2=-1$, $a^2=1$、
すなわち一鎖と二鎖の各一類で尽きる。全て半単純なBell直線の無印同型類と相異なる。
\end{theorem}
\begin{proof}[全ての四対応]
第一の交代成分の半単純化を前節の $P(s)$ とし、$E_{abcd}$ を二進語上の単位係数テンソルとする。
一鎖では固有値の順と共通倍率を選び、加える項を $E_{0011}$ とする。
二鎖ではこれに $E_{0110}$ または $E_{1001}$ も加える。それぞれ $+,-$ 型と呼ぶ。
実際、第一切断の二つの双曲平面は $00,11$ と $01,10$ の対である。
その各平面上でこの行列は対角 $1$ または $s$ に一つ以下の非対角行列単位を加えた形で、
右Gramの冪零部はその行列単位の非零倍となる。
自己随伴な二鎖の核は等方直線であり、双曲対の入替えと逆倍率でどちらの向きの非零係数も作れる。
二平面の向きを同時に反転することは特殊直交で、一鎖なら他平面で向きを合わせられる。
二鎖なら相対向き二通りを残す。右Gramの特殊直交同値と第一切断の行列式一致から
左側も特殊直交で同定される（定理\ref{thm:split-jordan-transport}）。

$P=P(s)+E_{0011}+\beta E_q$、$\beta=0,1$ とし、$\bar q$ を全ビット反転とする。
二逆転置の通常の成分は
\[
 I_0=P(1/s)-E_{1100}-\beta s^{-2} E_{\bar q},\qquad
 I_1=P(1/s)-s^{-2}E_{1100}-\beta E_{\bar q}.
\]
第一の四対応を恒等にしたとき、$Q=\pi_0 I_0$ に $\pi_1(\bigotimes_i L_i)I_1$ を射影的に合わせる。
比較は全て単項行列で行えるし、これで全てである。実際、差のある二固有値の一般固有平面の対を
右側および左側において保つか交換する必要がある。行列積の二側において各平面の二本の
可分解線が四つの二進座標線を定める。各逆行列の二つのブロックは対角または三角であり、
その可逆行列の台の完全マッチングは一つだけなので、四つの対角係数の台も保存される。
ここでは $\pi_1^{-1}\pi_0=(02)(13)$ が第一切断の二側を交換することを用いる。
従って必ず半単純部の比較にもなり、前の $T(c,d)$ と反転のリストを使える。

$F$ を伴う場合、第一の付加語を比較すると $cd=s^{-2}$、第二があれば
$+$ 型で $d/c=s^{-2}$、$-$ 型で $c/d=s^{-2}$ を得る。
$L=T(a,s^{-2}/a)F$ の全解が表の三行となる。反転なしでは第一の付加語の位置が合わない。
片側反転は $s^2=-1$ に限り、$s=\iota$, $\iota^2=-1$ を選ぶ。
一鎖ではその付加語が合わず、二鎖では $+$ 型で $L=T(e,\iota e)A$、
$-$ 型で $L=T(\iota e,e)B$, $e=\pm1$ のみが残る。
実際 $+$ 型の追加二式は $cd=\iota$, $c=-\iota d$、$-$ 型では
$cd=\iota$, $c=\iota d$ であり、反対の片側反転では台が合わない。
これで全表示を得た。正則性は定理\ref{thm:split-inverse-regularity}による。
\end{proof}

\begin{proof}[量子部分と追加二語による無印同型]
$F$ の行では両色のビットを位置2,3で反転する。六ブロックの非零語は24置換語と、
各辺 $uv\in E$ に対する $uvuv,vuvu$ だけである。置換語だけを残せば表の $T$ の量子窓となる。
一辺または共通端を持つ二辺なので、一つの共通端の文字だけを高重みにすれば量子窓へ流入する。

一鎖の場合、全位置に同じ対角行列を作用させ、$l_0+l_2=l_1+l_3$ とした二次元トーラスが
各位置で相異なる四指標を持つ。これは極大である。実際、分解トーラスの中心化群の階数は、
共通法線倍率の1に $P$ の安定群の階数を足した数以下である。
$P$ を保つ単項行列の対角部分は $T(c,d)$ で、$E_{0011}$ を保つため $cd=1$、
従って階数は高々1である。よって同型があれば極大トーラスの共役後に単項同型がある。
二鎖ならさらに $c/d=1$ が必要で、$P$ の安定群は有限、分解トーラス自体が極大である。
その共役後、同型は両側の二平面を保つか同時に交換する。後者でも旧 $Q$ は逆転置と
二側の交換から得られ同じ二固有平面と二鎖を持つ。各可分解線と第一の四対応から
全八つの二文字基底が線の順序・倍率まで定まるので、やはり単項同型で比較できる。

一辺の台では各位置の二文字の出現回数が7、他二文字が6で、追加語は7回の二文字だけを使う。
二辺の台では出現回数は $8,7,7,6$ で、追加語は6回の文字を含まない語として認識できる。
従って置換語の部分は保たれる。それを保つ位置別の置換は全位置で同じ置換である。
異なる二位置で相異なる入力文字を同じ出力にすることはできないからである。
これで $E,T$ の必要性を得る。

逆にその置換をした後、同じ $T$ の二つの置換語部分は位置別対角倍率で一致させられる。
これは三角積の等しい交換係数の比が文字倍率比となることと量子最上直線の一意性による。
相補二文字を $b,d$ とし、各辺の追加係数の比を測る次の量は、位置別の対角変更に不変である。
\[
   \frac{w_{uvuv} w_{vubd} w_{bdvu}}{w_{vuvu} w_{uvbd} w_{bduv}}=-1.
\]
相補文字の順序では量子切断内の階数1から値は変わらない。
\path{computations/verify_split_two_eigenvalue_jordan.py} はこの等式、$T$、および一巡回後も同じ値を
取ることを全行で有理関数として確認する。ゆえに各辺に残るのは両追加項共通の倍率だけである。
全位置に同じ対角基底を使うと、これは $d_ud_v/(d_bd_d)$ 倍になる。
二辺ならその二指標は独立なので、代数閉体上で二倍率を同時に所望の値に取れる。
これで十分性と、量子部分の一移動と同型なことを用いて最小周期1を得る。
\end{proof}

\begin{proof}[例外一類と成分]
二鎖の片側反転四表示を上の検算で相互に同定する。具体的には28個の非零座標で
位置別文字の出現回数を保つ置換だけを列挙する。対応後の係数比は全て四乗根である。
位置別の16個の倍率による作用の整数行列は語ごとにその四文字に1を置いた列を持つ。
Smith標準形で整数核の基底を取り、係数比の四乗根指数との積が全て法4で零なら
所望の倍率がある。四表示の各々と、代表の一移動についてこのような同型が得られる。
Jordan項は、$T(c,d)$ の適切な一径数で両項の重みを上げて消せる。
補完の半単純化は、二鎖の重みフィルトレーションによる定理\ref{thm:split-polystable-good-limits}
の操作である。ここでは四対応が単項で鎖フィルトレーションを運ぶので次数付きでも
同じ四対応が残る。より明示的には $+$ 型は $T(t^{-1},1)$ と旧第一対応で決まる相補作用、
$-$ 型は $T(1,t^{-1})$ を使い、両付加語は重み1となる。片側反転の四対応はこの変更で不変である。
この極限は半単純な片側反転の $a=\pm1$ の類であり、量子閉軌道と相異なる。
同じ分解において非半単純な $P,Q$ を持ち、階数1による分解の一意性から、
半単純直線の階数1のどの類とも同型にならない。一鎖の階数2では四指標が相異なり、
半単純直線の非量子階数2に現れる重複指標と区別できる。量子窓自身の階数は3以上である。

量子閉軌道を持つ行では定理\ref{thm:quantum-ore-incidence}により、十五次元成分に入る必要十分条件は
同じ成分の量子類を閉軌道に持つことである。四面の向きに合わせればこの条件は
$T=(\iota,-\iota,\iota,-\iota)$ またはその逆である。表の式に代入して $s^2=-1$, $a^2=1$ を得る。
この位数四の背景では偶置換が向き付けた面値を保ち、奇置換がその逆へ運ぶ。
無向一辺にも共通端を持つ無向二辺にも奇置換の安定化元があるので、各場合ちょうど一類である。
片側反転に流入する例外は、その閉軌道が十五次元成分に入らないため除かれる。
\end{proof}
